% sections/D_s2s.tex -- App. D, supplementary results of Sec. 4 (saddle-to-saddle). Material moved out of the
% main text in the final edit (2026-10-10, notes/report_style.md), with its labels: tables, secondary checks,
% per-seed numbers, the old Sec. 5.3 stall test, and the provenance and review history of Sec. 4.
\subsection{Saddle-to-saddle: supplementary results}\label{app:supp-s2s}

Details of \cref{sec:s2s}, in the order of its subsections. Tags as in the main text; settings in
\cref{app:exp-s2s,app:exp-locality,app:exp-order}.

\subsubsection{Phenomenology}\label{app:supp-s2s-phen}

\paragraph{Scope.}
\empirical{} Complete runs (all three components learned): binary4 320 of 320 (fan-in, both teachers, $\chi=3,6$) and
60 of 60 (balanced); MPS4 160/160 at $\chi=6$, 106/160 at $\chi=3$ and 60/60 balanced; binary8 50/50 and HT4 45/50 at
$\chi=6$ ($\isc=0.07$--$0.3$). At $\chi=3$ few runs with $\Kc=7$ complete (binary8 14 of 180, HT4 12 of 150), and no
MPS8 run learns its third component (0 of 284; 132 learn two; \cref{sec:s2s-failure}). The events of the two runs of
\cref{fig:s2s-runs} are at $3.6\cdot10^3$, $4.4\cdot10^4$, $5.6\cdot10^5$ (MPS4) and $1.3\cdot10^5$, $9.8\cdot10^5$,
$6.8\cdot10^6$ (binary8), with widths below $2$.

\paragraph{Plateau levels and spectra.}
\empirical{} Largest relative deviation of the loss from \eqref{eq:s2s-levels} at plateaus between resolved events
($\sep i>3$): $8.5\cdot10^{-3}$ (binary4, separated teacher, $\isc=0.1$), $2.0\cdot10^{-2}$ (binary4, close teacher,
at $\isc=0.003$ on a plateau with $\sep{}=5.1$), $2.8\cdot10^{-3}$ (MPS4; $8.1\cdot10^{-3}$ balanced), and
$2.1\cdot10^{-5}$ and $1.5\cdot10^{-4}$ at $\Kc=7,8$ ($\chi=6$, $3$); at $\isc=0.003$ at most $7.4\cdot10^{-7}$,
except on binary4 with the close teacher. In runs with the saddle-to-saddle signature the top $\abs S$ singular
values of every edge equal the learned weights to a relative $7.7\cdot10^{-3}$ (binary4, close), $9.0\cdot10^{-4}$
(binary4, separated; $1.7\cdot10^{-9}$ at $\isc=0.003$) and $3.5\cdot10^{-4}$ (other trees). The rank steps count
the singular values above $f\tw r$, $r$ the component being learned, with $f=\frac12$ or a nearby value that
separates $\tw r$ from the learned weights; the 34 exceptions among 531 runs are 28 binary4 runs with the close
teacher at $\isc\ge0.05$, whose events overlap, and 6 HT4 runs.

\paragraph{Separation of the events and the gap fits.}
\empirical{} Every separation ratio exceeds $3$ in all complete runs on MPS4 (326), binary8 (83) and HT4 (74), and
in 218 of 220 on binary4 with the separated teacher; with the close teacher on binary4 in 5 of 20 runs at
$\isc=0.1$, rising to 20 of 20 at $\isc=0.003$, as close weights make escapes overlap unless $\isc$ is small.
\Cref{tab:s2s-gaps} lists the per-seed fits; binary8 at $\chi=3$, over up to $1.1$ decades, gives $4.93\pm0.05$.
At the plateau midpoints the raw Hessian gives the same slopes as $\Hperp$ to $0.004$, and per seed the ratio of
$\abs{\lambda_{\min}}$ to the balanced-route value reaches $20$ at $\Kc=7,8$.

\begin{table}[tbp]
  \centering
  \caption{Plateau durations and saddle type, odeco teachers, fan-in initialization. Gaps: per-seed exponent
    $p$ of $\tev{i+1}-\tev i\propto\isc^{-p}$, mean $\pm$ s.e.\ over seeds (run sets: \cref{app:exp-s2s}),
    for $\chi=3$ and $6$ or, in rows ``both'', the separated and close teachers; brackets: range of the
    per-seed values (binary4, MPS4: over their four cells). $\lambda_{\min}$: mean per-seed slope of
    $\ln(-\lambda_{\min}(\Hperp))$ against $\ln\isc$ at the first and second plateau midpoints ($\chi=3$; $6$, or
    the two teachers). Predicted: $\Kc-2$. Ten seeds per fit except binary4 close (9, 9; 8, 7), MPS4 gap 2 at
    $\chi=3$ (6, 7), HT4 gap 2 (9), MPS8 (7; 3). $\isc$: $0.003$--$0.1$ (binary4, MPS4), $0.07$--$0.3$ at
    $\chi=6$, $0.03$--$0.4$ (binary8) and $0.1$ or $0.15$ to $0.4$ (MPS8) at $\chi=3$; $\dagger$: limited range
    (less than one decade, or fewer than five seeds; binary8 at $\chi=3$: $1.1$ decades for 11 of 20 seeds; some
    binary4 close-teacher seeds keep four or five values of $\isc$, about $0.8$ decade). Source:
    \texttt{experiments/s2s\_core/results/main\_exponents.csv}.}
  \label{tab:s2s-gaps}
  \footnotesize
  \setlength{\tabcolsep}{3pt}
  \begin{tabular}{@{}llc lll@{}}
    \toprule
    tree & teacher & $\chi$ & $\tev2-\tev1$ & $\tev3-\tev2$ & $\lambda_{\min}$ slopes \\
    \midrule
    binary4 & sep. & 3, 6 & $1.003\pm0.001$, $1.007\pm0.001$ & $1.004\pm0.003$, $1.011\pm0.002$ & 1.05, 1.02; 1.08, 1.04 \\
    ($\Kc=3$) & close & 3, 6 & $1.004\pm0.010$, $0.999\pm0.008$ & $1.016\pm0.011$, $1.014\pm0.006$ & 1.17, 1.12; 1.18, 1.17 \\
            & & & [0.95, 1.06] & [0.99, 1.09] & \\
    MPS4    & sep. & 3, 6 & $2.000\pm0.000$, $2.001\pm0.001$ & $1.990\pm0.005$, $1.998\pm0.001$ & 2.01, 1.99; 2.00, 2.00 \\
    ($\Kc=4$) & close & 3, 6 & $2.001\pm0.001$, $2.001\pm0.001$ & $2.004\pm0.002$, $2.001\pm0.001$ & 2.01, 2.00; 2.02, 2.00 \\
            & & & [2.00, 2.01] & [1.97, 2.01] & \\
    binary8 & sep. & 6$^\dagger$ & $5.020\pm0.009$ [4.98, 5.05] & $4.969\pm0.022$ [4.88, 5.09] & 5.10, 5.07 \\
    ($\Kc=7$) & both & 3 & $4.93\pm0.05$, $4.92\pm0.04$ [4.65, 5.13] & -- & 5.03; 5.03 \\
    HT4 ($\Kc=7$) & sep. & 6$^\dagger$ & $5.10\pm0.08$ [4.83, 5.78] & $4.75\pm0.14$ [3.93, 5.33] & 5.13, 4.88 \\
    MPS8    & sep. & 6$^\dagger$ & $5.95\pm0.05$ [5.79, 6.20] & not reached & 5.99 \\
    ($\Kc=8$) & both & 3$^\dagger$ & $6.0\pm0.3$, $6.2\pm0.6$ [5.14, 7.21] & not reached & 5.88; 6.07 \\
    \bottomrule
  \end{tabular}
\end{table}

\paragraph{The strict wrong-order saddle on binary4.}
\empirical{} With the close teacher on binary4 ($\chi=3$), seed 8 learns in the order 1-3-2 at seven of eight
values of $\isc$ ($0.005$--$0.1$), and seed 5 at two. At the plateau where $1$ and $0.7$ are learned and $0.85$ is
not, $\lambda_{\min}$ is $-0.106$ to $-0.146$ at the nine midpoints (raw: $-0.113$ to $-0.149$) and flat in $\isc$
(seed 8: slope $0.03$), against $-0.133$; the next event follows within about one width ($\sep{}=0.76$--$1.25$),
and this gap does not grow as $\isc\to0$ (seed 8: $\propto\isc^{0.08}$). At $\isc=0.003$ seed 8 learns greedily.

\subsubsection{Aligned initialization: integrator test}\label{app:supp-s2s-aligned}
\empirical{} Separated teacher, $d=4$, $\chi=3$, balanced aligned initialization with per-core amplitude $q_0=0.01$
(binary4), $0.07$ (MPS4), $0.29$ (binary8, HT4) and $0.34$ (MPS8), so that the events fall near $t=33$, $50$, $100$
on every tree. Gradient flow on the full TTN (Dormand--Prince 5(4), tolerances $10^{-10}$ and $10^{-14}$) follows
the exact solutions (\cref{fig:s2s-aligned}): the largest deviation of any overlap is $2.0\cdot10^{-10}$ (binary4),
$3.2\cdot10^{-10}$ (MPS4), $2.3\cdot10^{-9}$ (binary8), $2.6\cdot10^{-9}$ (HT4) and $2.7\cdot10^{-9}$ (MPS8); event
times agree on the same output grid to $1.3\cdot10^{-11}$ relative; the conservation drift \eqref{eq:drift} stays
below $3\cdot10^{-14}$. LSODA, the integrator of all other runs of \cref{sec:s2s}, deviates by at most
$7.1\cdot10^{-7}$. The test checks contractions, gradients, integrators and event detection on all five trees.

\begin{figure}[tbp]
  \centering
  \includegraphics{s2s/aligned.pdf}
  \caption{Aligned balanced initialization: overlaps from gradient flow on the full TTN (solid, Dormand--Prince,
    tolerances $10^{-10}$ and $10^{-14}$) and the exact mode solutions (dashed); separated teacher $(3,2,1)$, $d=4$,
    $\chi=3$; one run per tree, with per-core amplitude $q_0=0.01$ (binary4), $0.07$ (MPS4), $0.29$ (binary8, HT4)
    and $0.34$ (MPS8). Panel titles: largest deviation over time and components. Takeaway: with aligned
    initialization every tree learns the components in the order of their weights, at times $\propto1/\tw r$, as
    the decoupled modes of a deep linear network do (source: \cref{app:experiments}, \cref{tab:exp-s2s}).}
  \label{fig:s2s-aligned}
\end{figure}

\subsubsection{Escape from the origin}\label{app:supp-s2s-escape}

\paragraph{Test, seed by seed.}
\empirical{} The 231 pairs of initialization and teacher were run at eight values of $\isc$: $0.002$--$0.1$ on
binary4 and MPS4 ($\chi=3,6$, both teachers, fan-in and balanced) and $0.01$--$0.4$ at $\Kc=7,8$ (mostly separated
teacher, fan-in; \cref{tab:exp-s2s}). At the smallest $\isc$ the prediction \eqref{eq:s2s-first} matches every run to
a relative $8.5\cdot10^{-6}$ (binary4, separated teacher), $2.2\cdot10^{-3}$ (binary4, close teacher, whose direction
has not converged at $\isc=0.002$), $6\cdot10^{-5}$ (MPS4) and $4\cdot10^{-8}$ (binary8, HT4, MPS8 at $\isc=0.01$).
With the offset removed the per-seed exponents are $0.9996$--$0.9997$ ($\pm0.0001$; binary4, separated),
$0.990$--$0.996$ (binary4, close: $2.6$--$3.8$ s.e.\ below $1$, as an unconverged offset makes them),
$1.9999$--$2.0000$ (MPS4), $4.996$--$4.999$ (binary8, HT4) and $5.9995\pm0.0002$ (MPS8).

\paragraph{Dependence of the prefactor on $d$, $\chi$ and $\Kc$} (\cref{fig:s2s-prefactor}).
The leading-order flow gives $\mathcal{T}=\lim_{\isc\to0}\isc^{\Kc-2}\tev1$ per initialization cheaply, so these were measured
on many initializations ($d=4$, one component of weight $1$, unless stated). \emph{Physical dimension}
\empirical{} \newresult[the closed form]: at full bond $\chi=d^2$ (200 initializations per point) the binary4
medians are $1.544$ $[1.468,1.594]$, $1.569$ $[1.525,1.615]$ and $1.563$ $[1.525,1.587]$ (95\% bootstrap intervals)
for $d=2,3,4$, independent of $d$ and equal to $\pi/2$, the full-bond value of the closed form of
\cref{prop:A-product-mode} (\derived{} given end directions that do not rotate), while $t_{\rm KH}$ falls by a factor
$1.9$; $\mathrm{sd}\ln\mathcal{T}=0.218\pm0.010$, $0.157\pm0.009$, $0.135\pm0.006$ against the closed form's
$1/(2\sqrt\chi)=0.250$, $0.167$, $0.125$ ($\pm0.02$). MPS4 gives $1.39$ $[1.27,1.58]$, $1.41$ $[1.30,1.52]$ and $1.50$
$[1.43,1.60]$ (another set of initializations: $1.30$, $1.35$, $1.46$) while $t_{\rm KH}$ falls by $2.9$.
\emph{Bond dimension} \empirical{} \newresult: the median falls slowly, with log-log slopes $-0.335\pm0.028$
(binary4) and $-0.264\pm0.038$ (MPS4) over $\chi=4$--$16$ and $-0.76\pm0.08$ (binary8, $\chi\le12$), $-0.69\pm0.08$
(HT4) and $-0.72\pm0.06$ (MPS8) over $\chi=2$--$16$, against $-0.47$ and $-1.5$ to $-2.5$ for $t_{\rm KH}$ (another set
of initializations: $-0.20\pm0.05$ for MPS4, $-0.36$ to $-0.38$ for binary4). \emph{Number of cores} \conjecture{}
(heuristic) \empirical{} (agreement) \newresult: over 17 trees with $\Kc=3$--$15$ ($\chi=3$, weights $(3,2,1)$, 50
initializations each) the median ranges from $0.42$ (HT3, $\Kc=5$) to $814$ (binary16, $\Kc=15$), growing with $\Kc$
within each family (binary4 $0.86$, MPS4 $0.67$; at $\Kc=7$, HT4 $0.58$ and binary8 $3.04$). The RMS-overlap
heuristic $P_3=d^{\Ns(\Kc-2)/(2\Kc)}/((\Kc-2)\tw1)$, the escape time of a balanced aligned mode started at the
root-mean-square overlap $\isc^\Kc d^{-\Ns/2}$ of a fan-in $\Wz$ with a unit rank-one tensor, tracks it within a
factor $0.87$--$2.13$ with no free parameter, while the median of $\mathcal{T}/t_{\rm KH}$ grows from $3.1$ to $2.9\cdot10^7$.
\emph{Gauge balance} \empirical{} \newresult[minor]: re-gauging the same $\Wz$ to $\Q e=0$ changes $\mathcal{T}$ per seed by
a factor $0.88$--$1.92$ (cell medians $0.99$--$1.18$), keeps $r_1$ in 88 of 94 seeds and leaves the exponent
unchanged.

\begin{figure}[tbp]
  \centering
  \includegraphics{s2s/prefactor.pdf}
  \caption{Prefactor $\mathcal{T}=\lim_{\isc\to0}\isc^{\Kc-2}\tev1$ from the leading-order flow, medians over initializations
    with 95\% bootstrap intervals. Colour: tree family (MPS blue, binary orange, HT green); marker in (a), (b):
    circles $\Ns=4$, squares $\Ns=8$. (a)~Against $\chi$: one component of weight $1$, $d=4$, 100 initializations
    per point; dashed (open markers): median Kumar--Haupt time $t_{\rm KH}$. (b)~Against $d$ at full bond
    $\chi=d^2$, binary4 and MPS4, 200 initializations; dotted: $\pi/2$; dashed: $t_{\rm KH}$. (c)~Against $\Kc$, 17
    trees (diamonds; binary, MPS and HT with $\Ns=3$--$16$), $\chi=3$, $d=4$, weights $(3,2,1)$, 50 initializations;
    dotted: the RMS-overlap heuristic $P_3$; dashed: $t_{\rm KH}$. Takeaway: weak $\chi$-dependence, none on $d$ at
    full bond, $P_3$ follows the growth with $\Kc$, and $t_{\rm KH}$, which uses only the norm of the
    initialization, is low by factors that grow with $d$, $\chi$ and $\Kc$ (source: \cref{app:experiments},
    \cref{tab:exp-s2s}).}
  \label{fig:s2s-prefactor}
\end{figure}

\subsubsection{Locality: two components}\label{app:supp-s2s-locality}

\begin{table}[tbp]
  \centering
  \caption{Plateau law after the first component, two-component pattern teachers ($d=\chi=2$, weights
    $(1,0.5)$): prediction \eqref{eq:locality-rule} against measurement. $p$ is the exponent of
    $\Delta t_2\propto\isc^{-p}$, fitted per seed; for $k=2$ the slope of $\Delta t_2$ against $\ln(1/\isc)$.
    Ranges are over rows (one row = tree, $D$, protocol); single rows give mean $\pm$ s.e.\ over seeds (median).
    $\sep{}$ is the separation ratio \eqref{eq:separation}. $\dagger$: limited range (less than one decade of
    $\isc$, or fewer than five fitted seeds; the reporting rule of \cref{sec:setup-diagnostics}). Sets $D$ number
    the sites from $0$ (source: \cref{app:experiments}, \cref{tab:exp-locality}).}
  \label{tab:locality-exponents}
  \small
  \begin{tabularx}{\textwidth}{@{}c>{\raggedright\arraybackslash}p{3.0cm}
      >{\raggedright\arraybackslash}p{1.9cm}>{\raggedright\arraybackslash}X@{}}
    \toprule
    $k$ & rows & prediction & measured \\
    \midrule
    1 & binary4, binary8, star3 (4 cases) & no plateau & small: $\Delta t_2\propto\isc$, at most
      $8\cdot10^{-4}$, $\sep{}<5\cdot10^{-4}$; saddle: $\Delta t_2=0.63$--$0.65$, independent of $\isc$, $\sep{}=0.34$--$0.35$ \\
    2 & MPS, midroot6 (7 cases) & log, slope $1$ (small), $2$ (saddle) & $0.980$--$1.009$ (small),
      $2.002$--$2.071$ (saddle); log law preferred in 14 of 14 rows \\
    3 & all nine trees (22 cases) & $1$ & $0.98$--$1.04$ (42 rows) \\
    4 & MPS, star3, midroot6 (8 cases) & $2$ & small $1.95$--$2.05$ (MPS4 $\{0,3\}$: 3 seeds$^\dagger$), balanced
      $1.94$--$2.02$, $\chi=4$: $1.99$--$2.03$; saddle $1.48$--$2.39$ ($1.81$--$2.19$ on never-stalling seeds) \\
    5 & binary8, HT4 (6 cases) & $3$ & $2.89$--$3.13$, all protocols \\
    5 & HT8 $\{1,2\}$ & $3$ & saddle $2.74\pm0.37$ ($2.96$); $3.10\pm0.11$ on never-stalling seeds \\
    5 & MPS6 $\{1,3,5\}$, MPS8 $\{1,5\}$ & $3$ & $\isc\in[0.03,0.3]$: $2.86$--$3.76$.
      $\isc\in[0.012,0.03]^\dagger$: MPS6 $3.05\pm0.06$ (saddle), $3.15\pm0.23$ (balanced); MPS8
      $3.13\pm0.25$ (balanced, 4 seeds), $\mathbf{3.67\pm0.20}$ (saddle, 4 seeds). MPS8 small, never-stalling
      seeds: $2.97\pm0.05$. $\chi=4$: $3.19\pm0.15$, $2.97\pm0.04$ \\
    6 & binary8 $\{0,2,4\}$, HT4 $\{0,1,2\}$ & $4$ & balanced, one decade, 8 and 7 seeds:
      $4.007\pm0.021$, $4.061\pm0.048$ (fan-in$^\dagger$: $3.68\pm0.20$, $4.23\pm0.10$) \\
    6 & MPS8 $\{1,6\}$ & $4$ & balanced, one decade: $4.51\pm0.19$ \\
    7 & HT4, all sites & $5$ & balanced, one decade, 8 seeds: $4.957\pm0.061$ ($5.01$) \\
    6, 7 & other rows (6 and 3 cases) & $4$, $5$ & fan-in and saddle$^\dagger$, $1$--$5$ fitted seeds:
      $3.4$--$5.1$ ($k=6$), $5.3$--$6.3$ ($k=7$; HT8 $\{3,4\}$: $5.74\pm1.06$, median $5.01$) \\
    8 & MPS8 $\{0,7\}$ & $6$ & not measured ($\le2$ fitted seeds per protocol$^\dagger$) \\
    \bottomrule
  \end{tabularx}
\end{table}

\paragraph{Head start.}
\empirical{} MPS8 $\{3,4\}$ from small initialization nearly ties in the fair comparison (residuals $0.004$ and
$0.005$). On MPS4, $\tw2=0.25$ gives slopes $3.023\pm0.010$ (small) and $4.04\pm0.05$ (saddle) against $3$ and $4$,
and $\tw2=0.75$ gives $0.367\pm0.025$ and $1.345\pm0.016$ against $1/3$ and $4/3$. With weights $(2,1)$ from the
saddle on MPS4, MPS6, MPS8 (8 seeds each), $\lambda_{\min}=-1.4142$, $-1.5874$, $-1.6818$ against
$\lambda_2=2^{(\Kc-2)/\Kc}=1.4142$, $1.5874$, $1.6818$, and the slopes are $0.712\pm0.007$, $0.649\pm0.016$,
$0.595\pm0.003$ against $1/\lambda_2=0.707$, $0.630$, $0.595$.

\paragraph{Saddle type.}
\empirical{} Near rank-deficient states the gauge projection removes part of the escape direction
(\cref{app:exp-locality}), so raw values are quoted where the two differ; the fixed threshold of
\cref{sec:setup-diagnostics} would call $k=3$ plateaus strict for $\isc\gtrsim0.01$--$0.04$. For $k=2$ from the
saddle, $\lambda_{\min}=-0.5000$, projected and raw ($-0.2500$ and $-0.7500$ for $\tw2=0.25$, $0.75$). From small
initialization the raw value departs from $-\lambda_2$ where the plateau is too short to reach the saddle ($-0.437$
for MPS8 $\{3,4\}$ at $\isc=0.008$; $-0.581$ against $-0.75$ for $\tw2=0.75$). For $k\ge3$ the exponent $x$ of
$\lambda_{\min}\propto\isc^x$ ranges over rows from $0.81$ to $1.06$ for $k=3$ (raw $0.92$--$1.06$), $1.64$--$2.16$ for
$k=4$, $2.82$--$3.76$ for $k=5$, $4.00$ and $4.04$ for $k=6$ (balanced binary8 and HT4, one decade; $4.27$ on MPS8;
$3.19$--$5.23$ in the rows of limited range), and $5.01$ (balanced HT4) up to $6.5$ (limited range) for $k=7$; the
median of $\abs{x-p}$ is $0.02$; raw and projected exponents differ by up to $0.11$.

\begin{figure}[tbp]
  \centering
  \includegraphics[width=0.86\textwidth]{locality/twocomp_hessian.png}
  \caption{Saddle type from the Hessian at the plateau midpoint, same sweep as \cref{fig:locality-exponents}.
    Plotted: $-\lambda_{\min}$ of the gauge-projected Hessian $\Hperp$ (median over seeds) against $\isc$, one line
    per (tree, $D$), coloured by $k$; left small initialization, right saddle start; dotted:
    $\lambda_2=\tw2\tw1^{(\Kc-2)/\Kc}=0.5$. For $k=2$, $\lambda_{\min}=-\lambda_2$ at every $\isc$: a strict saddle
    (the small-initialization curves drop at the largest $\isc$, where the plateau is too short to reach the saddle).
    For $k\ge3$, $\lambda_{\min}\to0$ as $\isc\to0$, roughly as $\isc^{k-2}$: higher-order saddles. Lines for
    $k\ge6$ come from the rows of limited range of \cref{fig:locality-exponents}. Takeaway: with the larger component
    learned first, negative curvature that survives $\isc\to0$ exists only for two-core routes (source:
    \cref{app:experiments}, \cref{tab:exp-locality}).}
  \label{fig:locality-hessian}
\end{figure}

\paragraph{Bond spectra.}
\empirical{} Outside the Steiner subtree $\bs e2$ stays below $3\cdot10^{-3}\tw2$ for $\isc\le0.03$. The effective
rank uses the threshold $0.05$ times the teacher's value; with $0.03$ and $0.1$ the Steiner edges are the edges that
gain rank in $99.6\%$ and $99.9\%$ of the events. The spread of the crossing times of the Steiner edges has a 95th
percentile of $0.25$ event widths and grows as $\isc^{1.94}$.

\paragraph{Reading $k$ from a single run.}
On a balanced route the new overlap grows as $\dot{\ov2}\propto\ov2^{\,2-2/k}$ at small $\ov2$, so the local growth
exponent $\zeta=\dd\ln\dot{\ov2}/\dd\ln\ov2$ gives $k=2/(2-\zeta)$ (\cref{eq:A-growth-exponent}) \derived{}
\knownresult[the exponent $2-2/k$ of a balanced product of $k$ factors: \citealp{arora2019implicit,razin2022implicit}].
Fitted over $\ov2\in[10^{-4},10^{-2}]$, the medians per row give $k=2.00$, $3.0$, $4.0$, $5.1$--$5.3$, $6.0$,
$7.0$--$7.7$ for $k=2$ (saddle start) to $7$ \empirical{} \newresult[as a meter of $k$]. Validity range: from the
saddle, or from small initialization if the plateau is long enough for the window to lie after the escape
transient. The estimate fails for small-initialization $k=2$ on MPS6 and MPS8 ($2.4$--$3.3$; plateau shorter than
the transient), is biased high for $k\ge5$ (HT8: $5.3$ for $k=5$, $7.7$ for $k=7$), and cannot separate $k=7$ from
$8$, since $\dd k/\dd\zeta=k^2/2$ (per-run 5--95\% ranges $\pm1$ at $k=7$; MPS8 $\{0,7\}$, $k=8$: $7.9$ from small
initialization, 3 runs, $6.1$ from the saddle, 2 runs).

\paragraph{Independent cross-check.}
\empirical{} A pilot with separate code ($d=\chi=2$, weights $(1,0.5)$, MPS4, MPS6, binary4, binary8, fan-in
initialization; mean $\pm$ s.d.\ of per-seed fits over 3--8 seeds) gives the first escape $\tev1\propto\isc^{-(\Kc-2)}$
to within $0.003$ in the exponent for every seed; no plateau for $k=1$; a logarithmic $k=2$ plateau with slope
$1.01\pm0.01$ on MPS4 ($\isc=10^{-3}$--$0.1$; predicted $1$), rising to $2.3\pm0.4$ from the saddle (predicted $2$);
and $p=1.00$--$1.01$ for $k=3$, $2.00$ (MPS4) and $2.05\pm0.06$ (MPS6) for $k=4$, $3.00\pm0.11$ (binary8) for $k=5$
and $4.97\pm0.05$ (binary8, limited range: under one decade) for $k=7$. It also met the failure modes of
\cref{sec:s2s-failure}: 95 of 332 runs with $k\ge4$ (none with $k\le3$) never learn the second component, and on
MPS6 the exponents for $k=5$ and $6$ are seed-dependent and up to two units steeper than $k-2$ ($3.75\pm0.24$ for
$k=5$, with 3 of 6 seeds escaping).

\subsubsection{Several components: the saddle hopper}\label{app:supp-s2s-hopper}
Source: \path{notes/saddle_hopper.md} (v2) and \path{experiments/saddle_hopper/} (README there); audit
\path{notes/reviews/saddle_hopper_verification.md}, which finds the hopper reproducible but permissive. Section 9 of the
note (the theorem check) was added after the audit and is not covered by it.

\paragraph{What it is.}
At each plateau the hopper computes the fixed-rank (old) block and the local residuals of the connected core sets,
follows the most unstable direction, re-optimizes at the new bond-rank profile, and repeats; it adds the first escape
of \cref{sec:s2s-escape}, the race of the $k=2$ routes with the dormant amplitudes carried across saddles, the time
laws and the failure flags. It is deterministic (a re-run is byte-identical) and reproduces the worked examples of
the theory notes. Its predictions use neither $\isc$ nor the seed, and no rule is tuned per case; data-dependent
elements are one rule for the 1-3-2 order of binary4 (made explicit after reading the odeco runs), the saddle start
from the lowest rank-one attractor, and the checks of the causes of the mismatches, written after seeing them.

\paragraph{Scores} \empirical.
Of the 11,432 scored locality runs, 96.9\% follow the hopper's main branch (645 of them stalled runs admitted by its
trap flag, set on every route with $k\ge4$ under unbalanced initialization). On the informative runs (one main order,
not stalled; 7,965 runs) it matches 97.0\%, against 95.8\% for the weight order without rotation; on the fan-in and
balanced two-component runs the two predictions are identical (96.1\% each); on the 107 multi-component runs whose
main order is not the weight order (B4 on MPS4, E8 on MPS8) it matches 88.8\%, against 0\%. On the odeco runs it
matches 96.6\%, as does the weight order. Citable as a stall result: no stall at $k\le3$ (0 of 4,531 scored
two-component runs with $k=2$, $3$). Where it commits to one order, that order is the weight order in all but the 107
runs, so as a whole-sequence validation it is bookkeeping of known theory.

\paragraph{Multi-saddle slopes} \empirical{} \newresult.
For the nested teacher TC on MPS8 from the saddle, the race across saddles predicts the log slopes $1.667$, $5.356$,
$7.737$, then a $k=4$ plateau (exponent $2$), then $80.26$; measured: $1.718$, $5.540$, $7.581$, $2.018$, $78.92$ (all
five within $3.5\%$). From small initialization TC misses at its second level ($5.62$ against $7.74$), and TA and TB
are $10.5\%$ low at their first level (partial carry-over of the dormant amplitudes between saddles, not modelled).

\paragraph{The saddle-order theorem at the reachable saddles.}
\proven{} (landscape: \cref{thm:A-order-lower,thm:A-order-sharp} applied by exact computation) Every first-escape
candidate of the hopper gives a rank-one saddle, 587 in all, over 288 cases (tree, $d$, $\chi$, teacher). At 147 of
them $\lambda_{\max}(\Gamma)>\tw q$ ($q$ the learned component): strict saddles, wrong order or stars. At level
$\ell=2$ the generalized operator never exceeds $\tw q$; at level $\ell=3$ it does at 30 saddles, whose order is $6$
instead of $k^*=7$: locality HT4 ($\chi=2$, $4$) and HT8 with $D=\{0,1,2,3\}$, odeco runs on HT4 and binary8 with the
second or third component learned, and the binary8 runs of the order experiment (two components, $w=0.5$, $0.8$,
$0.95$; odeco with three and four components). All are reached after a wrong first escape; none is at a correct-order
saddle or on an MPS. \empirical{} The escape signature $x=(\Loss_S-\Loss)/(\tw r^2a)$, $a$ the overlap of the escaping
component $r$, is about $1$ for a route and $1-\tw q/\tw r$ for two symmetric clusters: in the odeco HT4 runs with the
close teacher, seed 4 gives $x=0.2998$, $0.2999$, $0.289$ at $\isc=0.07$, $0.15$, $0.2$ (prediction $0.300$) and a gap
exponent of $4.03$, against $\isc^{-5}$ for the route; most other wrong-first seeds take the route ($4.89$--$5.47$). These
escapes account for 19 of the 88 locality sequence mismatches, and for none of the steep MPS exponents and none of the
stalls. 108 later escapes of the hopper leave saddles with two or more learned components, which the theorem does not
cover.

\subsubsection{Order of learning}\label{app:supp-s2s-order}

\paragraph{The first component of the odeco runs.}
\empirical{} The one run of the first-escape test whose first component is not $r_1$ (binary4, $\chi=3$, close
teacher, balanced, seed 6, $\isc=0.1$) learns the weakest component first. On binary4 the largest component is
first in 639 of 640 first-escape runs and in all 380 full runs, which have 20 independent initialization
directions. On the tensor-like trees the largest component is first in 400 of 640 runs on MPS4, 264/328 on binary8,
104/160 on HT4 and 64/80 on MPS8; for a given seed it is the same at every $\isc$ (10 of 10 seeds in every cell of
the full runs). It is the LO flow's $r_1$ in all 231 runs at the smallest $\isc$. The race names it in all 640 MPS4
and 80 MPS8 runs, in 328 of 328 binary8 and 152 of 160 HT4 runs, and in 537 of 640 binary4 runs.

\paragraph{Tensor-like pairs.}
\empirical{} The race curves \eqref{eq:order-race} lie at $0.66$--$0.89$ ($w=0.5$) and $0.55$--$0.66$ ($w=0.8$), with
gradient flow $0.018\pm0.019$ below them over the 102 MPS cells, a fluctuation of the shared seeds. On binary8
$P(\text{smaller first})=0.240$, $0.4075$, $0.4875$ at $w=0.5$, $0.8$, $0.95$ for every $\isc\in[0.01,0.1]$; at
$\isc=0.3$ at most $1\%$ change.

\paragraph{Number of cores, bond and physical dimension; the selection rule.}
\empirical{} On MPS with $d=\chi=2$, $\mathrm{sd}\ln\hat t_1=0.59$, $0.93$, $1.20$ for $\Kc=4$, $6$, $8$. With
capped bonds the gain stops once the end bonds reach their cap $d$ (MPS8, $d=4$: $P(\text{larger first})=0.60$,
$0.73$, $0.76$ at $\chi=2$, $4$, $8$). With uncapped bonds $\mathrm{sd}\ln\hat t_1\propto\chi^{-(0.48\text{--}0.54)}$
($\chi\ge4$), and its growth with $\Kc$ approaches $\sqrt{\Kc-2}$. The test of $d$ at equal bonds (binary8 and MPS4,
independent seed sets) excludes differences above about $0.06$ at 95\%. With $X_r$ the projection of the
initialization onto the orbit of component $r$'s balanced rank-one configuration,
$\ln\hat t_{1,r}\approx\mathrm{const}-q\ln X_r$ with $q=(0.99$--$1.25)(\Kc-2)$ (s.e.\ $1$--$4\%$).

\paragraph{Matrix-like pairs.}
\empirical{} \Cref{tab:order-gamma} lists the fitted exponents. Fits over the whole measured range lie below
$(1-w)/w$ in 12 of 13 sets, by up to $24\%$: $2.5$--$3.3$ s.e.\ at $w=0.95$ and for MPS4 $D=\{0,2\}$ at $w=0.9$,
otherwise within $1.8$ s.e. For the adjacent pairs the fit over $\isc=10^{-5}$--$10^{-3}$ on MPS4 at $w=0.9$ gives
$0.40\pm0.05$; at $w=0.8$ and on MPS6 (down to $\isc=0.003$) the exponents stay below $n(1-w)/w$. With two components
binary4 is far greedier than a race of independent clocks ($0.87$ against $0.73$ at $w=0.8$).

\begin{table}[tbp]
  \centering
  \caption{Matrix-like pairs: exponent $\gamma$ of $P(\text{smaller first})=C\isc^{\gamma}$, binomial maximum likelihood
    over the whole $\isc$ range, s.e.\ from a bootstrap over seeds (shared across $\isc$); $d=2$, $\chi=4$ for the pattern
    pairs; weights $(1,w)$. Prediction \eqref{eq:A-matrix-like}: $n(1-w)/w$, $n=1$ for non-adjacent cores, $n=\bd e=4$
    (conjectured) for adjacent ones. --: not run, or (MPS6 $\{2,3\}$, $w=0.7$) events at two values of $\isc$ only.
    Source: \texttt{experiments/order/results/sigma\_fits.csv}.}
  \label{tab:order-gamma}
  \footnotesize
  \setlength{\tabcolsep}{4pt}
  \begin{tabular}{@{}llcccc@{}}
    \toprule
    pair & $\isc$ range, seeds & $w=0.7$ & $0.8$ & $0.9$ & $0.95$ \\
    \midrule
    binary4, $\chi=2$ & $10^{-4}$--$0.1$, 1,000 & $0.370\pm0.034$ & $0.237\pm0.018$ & $0.103\pm0.008$ & $0.042\pm0.004$ \\
    binary4, $\chi=4$ & $10^{-4}$--$0.1$, 1,000 & $0.42\pm0.05$ & $0.248\pm0.018$ & $0.102\pm0.008$ & $0.040\pm0.004$ \\
    MPS6, $D=\{0,5\}$ & $0.003$--$0.1$, 400 & $0.46\pm0.06$ & $0.249\pm0.029$ & $0.101\pm0.014$ & -- \\
    MPS4, $D=\{0,2\}$ & $10^{-5}$--$0.1$, 1,000 & -- & $0.240\pm0.017$ & $0.095\pm0.006$ & -- \\
    \multicolumn{2}{@{}l}{prediction $(1-w)/w$ ($n=1$)} & $0.429$ & $0.250$ & $0.111$ & $0.053$ \\
    \midrule
    MPS6, $D=\{2,3\}$ (adjacent) & $0.003$--$0.1$, 400 & -- & $0.67\pm0.09$ & $0.22\pm0.03$ & -- \\
    MPS4, $D=\{1,2\}$ (adjacent) & $10^{-5}$--$0.1$, 1,000 & -- & $0.82\pm0.06$ & $0.31\pm0.02$ & -- \\
    \multicolumn{2}{@{}l}{conjecture $4(1-w)/w$ ($n=\bd e=4$)} & $1.71$ & $1.00$ & $0.444$ & -- \\
    \bottomrule
  \end{tabular}
\end{table}

\paragraph{After a wrong first escape.}
\empirical{} Saddle starts: binary4 and MPS4 $D=\{0,2\}$ rotate (192 and 144 runs), MPS4 $D=\{1,2\}$ adds (144 runs,
$\ov2$ stays at $1$), the odeco MPS4, MPS6 and binary8 keep the wrong order ($\ov2\ge0.96$). At $w=0.8$ the midpoint
$\lambda_{\min}$ (raw) is $-0.18559$, $-0.17890$ and $-0.89427$ (three times) against $-0.18566$, $-0.17888$ and
$-0.89443$, independent of $\isc$; the escape slopes against $\ln(1/\isc)$ are $5.466\pm0.082$, $5.5903\pm0.0045$,
$1.1171\pm0.0009$ against $5.386$, $5.590$, $1.118$ (mean $\pm$ s.e.\ over seeds; the binary4 excess comes from
$\isc=10^{-2}$, and the local slope times $\abs{\lambda_{\min}}$ tends to $1.0000$ at $\isc=10^{-7}$), and a power law
fits $11$--$88$ times worse in rms residual ($95$--$106$ on the adding branch). On the kept branch the escape times
scale as $\tev1\propto\isc^{-1.98\pm0.03}$, $\isc^{-4.26\pm0.13}$, $\isc^{-5.3\pm0.3}$ on MPS4, MPS6, binary8 (MPS6 and
binary8 limited range); one MPS6 seed of six stalls. Natural runs (up to 40 per $(w,\isc)$): the scaled correction
time $\abs{\lambda_{\min}}(\tev1-\tev2)$ has slope $-0.02\pm0.03$ per unit of $\ln(1/\isc)$, and a wrong plateau is
resolved in at most $2.5\%$ of the runs at $\isc=10^{-4}$--$10^{-2}$; the smaller component overshoots
and falls to zero while the larger rises, then come a plateau at $\frac12w^2$ and the relearning. At $\isc=0.1$, 17,
12 and 1 of 40 binary4 runs rotate at $w=0.7$, $0.8$, $0.9$ (52 of 120 are in another class), with scaled times
$0.85$, $1.16$, $0.38$, and $5$--$7.5\%$ show a wrong plateau. All 150 adjacent runs add, within $0.1$--$1$ time units; on
MPS4 odeco the wrong plateau is resolved and 39 of 40 runs add after it, while one stalls.

\paragraph{Three and four components.}
\empirical{} One $\isc$ per tree ($0.05$ binary4, $0.1$ MPS4, $0.2$ MPS6, $0.25$ binary8), keeping dormant components
above the float64 floor; the last two lie outside the checked range of $\isc$-independence. MPS4: $0.35$ and $0.23$
greedy against $0.36$ and $0.25$ for the model. Paired with each run's own LO race, 41 runs are greedy only in
gradient flow and 43 only in the race; the own race names the first event in $97$--$100\%$ of the runs. Runs that stalled or did not finish by $t=10^9$ are excluded from both sides:
at most 5 of 200 on MPS4 and binary8, about half on MPS6. With four components and $\chi=4$, MPS6 stops before the
fourth in 396 of 399 runs and binary8 in 389 of 393, and none of 22 runs without any stall stop learns all four by
$t=10^9$.

\subsubsection{Failure modes}\label{app:supp-s2s-failure}
In the odeco runs the classes of \cref{app:exp-s2s} called \emph{trapped} and \emph{stalled} are the dead routes and
the slow routes of \cref{sec:s2s-failure}.

\paragraph{Odeco teachers.}
\empirical{} With the close teacher on MPS4 at $\chi=3$ the outcome is not fixed by the direction (one seed has a dead
route at 6 of 8 values of $\isc$). On HT4 at $\chi=6$ the 5 dead routes of 50 runs come from one seed, at every
$\isc$. The classes are unchanged at a 100 times smaller absolute tolerance (9 of 9 runs). The optimizer experiment
met the same failure on its odeco MPS8: fan-in gradient flow converged in none of 12 runs, with dead route ends in 11
(\cref{app:exp-optimizers}; a fix-round result of that experiment, not re-verified independently). Balanced initialization gives no dead route where fan-in gives 12 of 30 and 47 of 90; 17
of those 47 fan-in runs complete when re-gauged, and 5 runs that complete from fan-in initialization are slow when
re-gauged. The routes die on the plateaus, not in the escapes: across an escape a dormant amplitude
$\abs{\ip{\W}{\U r}}$ changes by a median factor $2.1$--$3.0$ at $\Kc=3$ (it grows: the head start of
\cref{prop:A-head-start}), $0.64$--$0.74$ at $\Kc=4$ and $0.12$--$0.75$ at $\Kc=7,8$, and falls by more than $100$ times
in at most $1.3\%$ ($\Kc=4$), $2\%$ ($\Kc=7,8$, $\chi=6$) and $8.3\%$ ($\chi=3$; $11\%$ for the later escapes on MPS8
with the separated teacher) of the escapes.

\paragraph{Slow routes: slow where tested.}
\empirical{} (prediction \derived) \newresult{} A slow route with live ends should escape after about the plateau of a
balanced route at its present amplitude, $b^{-(\Kc-2)}/((\Kc-2)\tw r)$ with $b=(\tw r\abs{\ov r})^{1/\Kc}$. Continued
from their final state, 16 of 23 slow or unresolved runs escaped, at $0.24$--$2.19$ times this estimate (all 8 with the
shortest estimates, 6 of 7 further runs of a seeded random sample, 2 of 8 MPS8 runs at $\chi=6$); 2 did not escape
within three times it and 5 ran out of wall-clock time or failed. The estimate is not a bound (4 escapes came later).
For the slow runs not continued (91 of 106 at $\Kc=7,8$ with $\chi=3$, 35 of 43 MPS8 runs at $\chi=6$), ``a long
plateau, not a dead route'' is a \conjecture: their dormant slices are pinned at $0.01$--$0.04\,\isc$ on binary8 and
balanced MPS8, but were still decaying in 6 of 13 balanced HT4 runs and 17 of 59 fan-in MPS8 runs.

\paragraph{Pattern teachers.}
\empirical{} A stalled two-component run sits at the partial fit $\Loss=\tw2^2/2$. Stalls by $k$, both protocols pooled,
fan-in initialization at $\chi=2$: $130$ of $768$ runs ($k=4$),
$62/816$ ($k=5$), $133/360$ ($6$), $29/118$ ($7$), $44/60$ ($8$). The dead-route check reran 22 (case, $\isc$) pairs
with a 60 times longer patience (170 runs); the 37 flagged runs end with route-end amplitudes of
$10^{-17}$--$10^{-13}$. The pilot, with separate code, found the same collapse to the rank-one saddle. Of 10 balanced
stalls rerun with longer patience, 6 escaped very late (after $2\cdot10^9$--$5\cdot10^{10}$, about $10^5$ times the usual
plateau), 3 are pinned (route ends at $3\cdot10^{-10}$, $6\cdot10^{-9}$, $4\cdot10^{-8}$ at $t=10^{11}$--$5\cdot10^{11}$)
and one was in the middle of a late escape when its budget ended; all 16 balanced stalls are on MPS routes. Spare
channels: on routes with an uncapped end the stalls fall from 79 at $\chi=2$ to none at $\chi=4$ (all 13 such cases
escape); on routes with both ends capped they rise from 90 to 106, in 7 of the 8 such cases that stall at all (they
fall only for MPS4 $\{0,3\}$, from 30 to 6 of 48). midroot6 and MPS6, the same graph with different roots, share
exponents but not stall rates ($D=\{0,5\}$ at $\chi=4$: 30 and 18 of 30): the root changes the fan-in statistics of
$\Q e$. The odeco runs agree: $\chi=6$ removes the dead routes on binary8, where every bond has twice the teacher
rank, but not on HT4 and MPS8, whose route ends sit on bonds capped at $d=4$, one channel above it.

\begin{table}[tbp]
  \centering
  \caption{Stalled runs among two-component runs with $k\ge4$: small initialization, $d=2$, the same 24 (tree, $D$)
    cases and $\isc$ grids in every column. Fan-in initialization at $\chi=2$ and $\chi=4$; balanced initialization
    ($\Q e=0$, same $\Wz$) at $\chi=2$. Bottom rows: split by whether at least one end of the route sits on a bond
    below the cap \eqref{eq:bond-cap} or both ends sit on bonds capped at $d$ (MPS chain ends, HT leaves). Source:
    \texttt{experiments/locality/results/stall\_capped\_ends.csv} and \texttt{run\_classes\_by\_sweep.csv}.}
  \label{tab:locality-stalls}
  \small
  \begin{tabular}{@{}lccc@{}}
    \toprule
    $k$ & fan-in, $\chi=2$ & balanced, $\chi=2$ & fan-in, $\chi=4$ \\
    \midrule
    4 & 67/384 & 2/384 & 18/384 \\
    5 & 20/384 & 5/384 & 0/384 \\
    6 & 58/180 & 5/180 & 59/180 \\
    7 & 8/30 & 0/30 & 12/30 \\
    8 & 16/30 & 4/30 & 17/27 \\
    \midrule
    all & 169/1,008 & 16/1,008 & 106/1,005 \\
    \quad route with an uncapped end (13 cases) & 79 & -- & 0 \\
    \quad both route ends capped (11 cases) & 90 & -- & 106 \\
    \bottomrule
  \end{tabular}
\end{table}

\paragraph{Effect on the exponents.}
\empirical{} Stalls depend on $\isc$: in 20 of 33 fan-in rows with stalls a stalling seed escapes at other values of
$\isc$, so it enters the fits with fewer points. Restricting each row to seeds that never stall leaves the exponent
unchanged in 27 of 46 rows with stalls and moves it by $0.1$ or more in 9 fan-in rows (by up to $1.38$; four of them
keep only 1--3 seeds) and 3 other rows. This accounts for the off-target saddle-start $k=4$ and HT8 rows of
\cref{tab:locality-exponents}.

\paragraph{The earlier reduced-model criterion.}
\empirical{} Evaluated once, at the start of the plateau ($\ov1=0.99$), in the 122 fan-in and saddle-start runs of the
stall diagnosis, an earlier form of the criterion of \cref{prop:A-trap}, which kept only the route entry of the
partner's block $Y_{\rm nv}$ (an end is unfavourable if $c=Q_{\nu\nu}+\norm{y_{\rm route}}^2<0$; dead if two or more ends
are), is strongly associated with death (30 of 55 runs with both
ends unfavourable died, 7 of 67 with a favourable end) but is neither sufficient (25 predicted-dead runs escaped) nor
necessary (seven dead runs had a live end at the plateau start). It conflicts with the data in three places:
(i)~exactly balanced initialization leaves no dead route in either experiment (consistent), but on the pattern-teacher
MPS8 routes the ends shrink to $10^{-10}$--$10^{-8}$ and escapes are delayed by $10^5$ or more (the balanced odeco MPS8
runs are slow instead, with route slices near $0.03\,\isc$); (ii)~its heuristic frequencies (about half of the
two-ended HT routes dead under fan-in initialization; fewer stalls on MPS with $\chi=d$) are contradicted: 0 of 144
small-initialization HT4 runs with $k=5$ stall, against 30 of 48 on MPS4 $\{0,3\}$ and 17 of 48 on MPS6 $\{0,3\}$ (the
odeco HT4 routes, with four leaf ends, end dead in at least 24 of 30 runs at $\chi=3$); (iii)~spare channels remove
stalls only on routes with an uncapped end, here as in the regime maps, and the stalls rise on capped routes. The
full criterion (\cref{app:supp-s2s-trap}) resolves (ii) and the drop in (iii).

\paragraph{Steep MPS routes.}
\empirical{} (mechanism open) One row stays steep at small $\isc$: MPS8 $\{1,5\}$ ($k=5$) from the saddle,
$p=3.67\pm0.20$ at $\isc\in[0.012,0.03]$ (median $3.83$; limited range: 4 fitted seeds, $0.4$ decade; 10 of 24 runs
stalled), with a matching Hessian exponent ($3.74$ projected, $3.75$ raw). Its fitted seeds escape at every $\isc$, so
it meets a refutation criterion of the theory (a per-seed exponent off by more than $0.2$ at $\isc\le0.05$ in runs
without stall or sign correction), on thin evidence. In the steep MPS rows the route-end amplitudes at the plateau
midpoint scale as $\isc^{a}$ with $a=1.24$--$1.91$ and shrink during the plateau (median ratio of midpoint to first
event $0.53$--$0.78$, against $0.85$--$1.03$ in on-target rows); but $a>1$ also occurs in on-target rows (binary8
$\{0,7\}$, balanced: $a=1.34$, $p=2.99$), so the steep exponents are associated with shrinking route ends, not shown
to be caused by them. Neither the logarithmic correction expected if small ends caught up through the conservation
laws (per seed it fits much worse than a free power law: median residual of $\ln\Delta t_2$ $0.25$--$0.62$ against
$0.05$--$0.12$) nor a law $\isc^{-(2k-6)}$ (both ends pinned, $a=2$; measured $a$ is seed-dependent, between 1 and 2)
describes the data.

\paragraph{Stall test of the regime maps.}
\empirical{} The test of \cref{tab:stall-test} ran budgets of $6\cdot10^4$--$1.6\cdot10^7$ time units. Fan-in
initialization makes the root bonds child-heavy, $\E\Q e=\isc^2(1-1/\chi)\Id\succ0$, so with $\chi=\ncomp$ the route
ends of the last component decay at a rate $\propto\isc^2$ while its escape needs a time $\propto\isc^{-(\Kc-2)}$. From
fan-in initialization the stalled route is dead: its third bond singular value ends below $10^{-12}$ in 28 of 32
stalled plain runs. With balanced initialization the route still decays and 12 of 18 runs do not escape: the third
singular value is $1.4\cdot10^{-5}$ (median) at the end. In the plain-model $\isc\times\Kc$ sweeps the stalled runs are
typed by the $\ncomp$-th bond singular value at the end (thresholds on the end value, not trends): below
$10^{-12}$, a dead route (546 of 795), between $10^{-12}$ and $10^{-4}$ (212) or above (37); 182 of the 795 stopped
at the step limit, and some of the 37 at $\isc\ge0.77$ learned nothing. In the stall test 9 of 18 fan-in MPS8 runs
learn only the first component. At $d=2$ the
MPS with full bonds still stall in up to 3 of 6 seeds. Some runs of the maps stop at the integrator's step limit:
unfinished lazy fits at square bonds (\cref{sec:other-lazy}) and some deep-tree runs, for example on MPS with $\Ns=8$
at $\isc=0.1$ (after $2\cdot10^4$ steps, at $t=1.4\cdot10^9$ of a budget of $5.7\cdot10^9$).

\begin{table}[tbp]
  \centering
  \caption{Stall test at $\chi$ equal to the teacher rank: $d=3$, $\chi=3$, odeco teacher $(3,2,1)$,
    $\isc\in\{0.3,0.45,0.6\}$, six seeds; ``balanced'': the same $\Wz$ with $\Q e=0$ on every bond. Types of the stalled
    runs from the third bond singular value of $\W$ at the end, minimized over edges: dead route below $10^{-12}$,
    decaying up to $10^{-4}$ (no stalled run is above); the centred runs did not record the spectra of $\W$. Source:
    \texttt{experiments/regime\_map/results/report/stall\_test\_summary.csv}.}
  \label{tab:stall-test}
  \footnotesize
  \begin{tabular}{@{}lllcccc@{}}
    \toprule
    model & tree ($\Kc$) & initialization & converged & stalled: dead route & decaying & not typed \\
    \midrule
    plain   & binary8 (7) & fan-in   & 4/18 & 12 & 2  & -- \\
            &             & balanced & 6/18 & 0  & 12 & -- \\
            & MPS8 (8)    & fan-in   & 0/18 & 16 & 2  & -- \\
    centred & binary8 (7) & fan-in   & 4/18 & -- & -- & 14 \\
            &             & balanced & 8/18 & -- & -- & 10 \\
            & MPS8 (8)    & fan-in   & 0/18 & -- & -- & 18 \\
            &             & balanced & 3/12 & -- & -- & 9 \\
    \bottomrule
  \end{tabular}
\end{table}

\subsubsection{Trap criterion from the initialization}\label{app:supp-s2s-trap}
Source: \path{notes/theory_trap.md} (v1) and \path{experiments/theory_checks/trap/} (README there), verified in
\path{notes/reviews/trap_verification.md}: tables reproduced, verdicts recomputed with independent code for a stratified
sample of 96 runs with $k\ge4$ (94 agree), the conservation law checked in full gradient flow; its errata are applied
here and in \cref{prop:A-trap}. Setting: a learned rank-one saddle of component $q$ (every core
$a\,p_v\otimes\bigotimes_fw_f$, $a=\tw q^{1/\Kc}$), the dormant component $r$ with route $C$, $k=\abs C$; an \emph{end} $u$
of $C$ has one route neighbour, its \emph{partner} $w$, on the route edge $e$.

\paragraph{Statements.}
(1)~On the plateau the entries with all bond indices learned relax fast; the entries with one new bond index form
two-core products whose Gauss--Newton cost is only their part normal to the learned manifold (lemma \proven; reduced
flow \derived). (2)~In this reduced flow $C_u=X_u^{\perp\top}X_u^\perp-X_w^{\perp\top}X_w^\perp$ is conserved, and the
end's route entry survives iff $C_u$ has a positive eigenvalue (\proven{} within the reduced flow; in full gradient
flow $C_u$ stays constant to four digits after the initial transient, and the dying blocks go to $10^{-30}$, in the
verification's run on MPS6). (3)~Identity: $C_u=\nu\T Q^{(u)}_e\nu+\norm{Y_{\rm nv}}^2-\norm{X_{\rm nv}}^2+O(\isc^3)$, the
conserved bond imbalance (end minus partner) plus the partner's non-visible new entries minus the end's; an earlier
form kept only the route part of $Y_{\rm nv}$ (exact identity; checked to 3--4 digits). (4)~A partner without sites
gives no second-order cost: the end then decays through a three-core cluster at a rate $\propto\isc^4$; in general an
end of visible degree $j$ decays at a rate $\propto\isc^{2(j-1)}$, so it can kill its route only if $k\ge2j$: $k\ge4$ for
a partner with sites, $k\ge6$ at HT and binary leaves, never for $k\le3$ (\derived). (5)~For $k\le5$, as $\isc\to0$, a
route is dead iff two or more of its ends die, and one dying end is pinned at an amplitude $\propto\isc^{k-3}$
(\derived); at $k=4$ a $\isc$-free competition, settled by a reduced route model, decides (\derived{} where every
route edge has one new direction; with $\chi\ge3$ on an inner edge the model's scalar route amplitude is a choice of
basis, and its verdicts change with it). From $k=6$ on, three-core clusters through inner route cores decay at a rate
$\propto\isc^4$ and are a further death channel, marginal at $k=6$ for a live route and dominant beyond, so the
two-end rule is incomplete there. That a pinned route escapes after $\propto\isc^{-(2k-6)}$ is a \conjecture: untested
at $k=5$, and at $k=6$ the verification's gradient-flow run (MPS6 $\{0,5\}$, seed 3, $\isc=0.03$) pins the dying end
and then decays instead of escaping. (6)~An end with more new directions on its route edge than its partner has
normal physical rows generically never dies (spare channels protect only there; \derived). (7)~The plateau-start
state is $\isc$ times a function of the initialization direction (checked on MPS4 $\{0,3\}$: $\isc=0.02$ and $0.005$
agree to 3 digits), so the outcome is a property of the direction (\derived{} as $\isc\to0$, with $O(\isc)$
corrections). (8)~Balanced initialization: $C_u\ge0$ at ends that are tree leaves, which cannot die at leading order;
ends with extra legs (pattern teachers on MPS) still can (\derived). No balanced odeco run ends with a dead route (0
of 150), but 69 of them stall (binary8 27, HT4 12, MPS8 30 of 30 each) \empirical.

\paragraph{Per-run tests} \empirical.
Every committed two-component run of the locality sweeps and of the pilot was rebuilt from its configuration and seed
and predicted without gradient flow. Fan-in initialization, saddle start and small initialization (3,492 locality
runs with $k\ge4$, 587 stalled; the 191 wrong-order runs excluded, which behave like the rest), stalled over runs per
predicted class: dead $408/457$ ($89\%$), pinned $141/934$ ($15\%$), alive $38/2{,}101$ ($1.8\%$); the earlier criterion
(route part of $Y_{\rm nv}$ only; dead if two or more ends are unfavourable) gives $417/1{,}143$ ($36\%$), $119/1{,}116$ and
$51/1{,}233$ ($4.1\%$). By $k$ (dead, pinned, alive): $k=4$: $203/232$, $2/432$, $5/796$; $k=5$: $73/79$, $10/116$,
$9/1{,}089$; $k=6$: $86/98$, $87/284$, $18/150$; $k=7$: $25/27$, $8/60$, $0/42$; $k=8$: $21/21$, $34/42$, $6/24$. Stall
diagnosis (60 times the patience): predicted dead 30 dead, 6 escaped, 2 pinned; predicted pinned 5, 63, 0; predicted
alive 2, 56, 0. Pilot, by its own collapse detector: 56 of 63 predicted-dead runs collapse, none of 150 predicted
alive. Counting also the pilot runs that reach the time or wall-clock limit with the second component below
$10^{-20}$ as dead (a rule added in this analysis; 46 runs, 10 of which end at $10^{-9}$--$10^{-4}$, not at roundoff)
gives $102/109$ predicted-dead runs dead, $20/182$ pinned and $2/146$ alive. MPS4 at $\chi=2$, fan-in (locality and
pilot): all 150 predicted-dead runs stall or die, none of the other 344. Balanced initialization: 16 of the 50
predicted-dead runs stall, none of them dead (very late escapes instead); 4 of 990 predicted-alive runs stall. Per seed
(cases with at least three values of $\isc$): 447 of 462 predicted-alive seeds never stall, 45 of 65 predicted-dead
seeds stall at every $\isc$ (16 at some). At $\Kc=8$ the correlation flow, stopped at the norm ratio $10^{10}$, leaves
the split into learned and small entries ambiguous at about $6\%$, enough to flip verdicts with $\abs{C_u}<0.003$ (2 of
the 96 recomputed). Known failures (56 false alarms and 40 misses among 3,929 fan-in runs with the pilot): $k=4$ with spare inner
channels (MPS4 at $\chi=4$: 20 of 46 predicted-dead runs escape; these verdicts depend on the basis, see (5)); the
marginal case $k=6$ on HT4 $\{0,1,2\}$ (11 of 18 predicted-dead runs escape); 24 predicted-alive runs on MPS6 and
midroot6 routes with $k=6$ that stall at the small end of a $\isc$ grid starting at $0.08$ (seed-specific, barely live
ends such as $C_u=+0.007$, or the inner clusters; not large $\isc$); routes with one dying end that end dead (5 runs of
the stall diagnosis, all at $\isc\ge0.19$; 20 pilot runs, the MPS6 $\{0,5\}$ family with seed 3, at every $\isc$): the
inner-cluster channel, not large $\isc$; at large $\isc$, two dead HT4 $\{1,2\}$ saddle-start runs at $\isc=0.19$,
$0.30$ predicted alive (the same seed stalls at all six $\isc$, as slow escapes) and three predicted-dead runs that
escape at $\isc\ge0.45$ (HT8 $\{0,1,2,3\}$ twice, MPS8 $\{1,6\}$); and seeds near the boundary at $k=4$--$5$. The stall rate of
pinned routes does not grow monotonically with $k$ ($31\%$, $13\%$, $81\%$ at $k=6$, $7$, $8$).

\paragraph{Odeco teachers.}
\empirical{} The dying component is the last one: the trap acts on the third plateau, after two escapes. The learning
order is predicted from the initialization (the first component in 1,196 of 1,196 runs, against 844 for ``largest
first''; the first two in order in 872 of 1,029), but the trap is not: from the initialization alone the criterion
does not separate the classes. From the state after the second escape (that is, with gradient flow up to the second
escape) it does, nearly: MPS4 at $\chi=3$, $\isc=0.03$, 10 of 10 seeds with the separated teacher and 9 of 10 with the
close one (the tenth, predicted to escape through an end at $C_u=+0.0097$, stalls at every $\isc$); HT4 at $\chi=3$,
$\isc=0.2$, 6 of 6 on seeds 0--5 and 8 of 10 with seeds 6--9 added (two runs with one dying end trap instead of being
pinned). A $\isc$-free model of the second plateau would close the gap.

\paragraph{Probability under fan-in initialization.}
\proven{} for one new direction on the end's route edge (saddle start, partner with sites, independent Gaussian
entries; the model is exact) With $n_u=d^{s_u}-1$ normal physical rows of variance $\isc^2/F_u$ at the end and $n_w$ of
variance $\isc^2/F_w$ at the partner ($F$ the fan-in, $s$ the number of sites), $p_u=P(C_u<0)=I_{F_u/(F_u+F_w)}(n_u/2,n_w/2)$:
$0.392$ at an MPS end that is a tree leaf ($d=\chi_e=2$, any $\chi$), $0.5$ at an MPS root end at $\chi=2$ ($0.392$ at
$\chi=4$); an internal MPS end has $0.5$ at $\chi=2$ and is protected at $\chi\ge3$ (Monte Carlo with $2\cdot10^5$ draws:
within $0.002$). The formula needs one new direction: at an end bond capped at $d\ge3$ there are $d-1$, death needs
$C_u\le0$ on all of them, and the probability falls to about $0.023$ ($d=3$) and $1.5\cdot10^{-4}$ ($d=4$) (Monte
Carlo). For a route with two ends ($k\ge5$, by the two-end rule of (5)): $P(\text{dead})=p_up_{u'}$, $P(\text{pinned})=p_u+p_{u'}-2p_up_{u'}$ ($0.196$ and $0.50$ for an
MPS chain with a leaf and a root end). A Monte Carlo with fresh seeds agrees with the closed form within two standard
errors at every case with $k\ge5$ \empirical; the trap probability is $0.1$--$0.3$ per run for routes whose two ends
have partners with sites (a four-ended HT4 route: $0.33$--$0.42$). The committed seeds of MPS4 $\{0,3\}$ (5 of 8
trapping) are a trap-prone draw (probability about $3\%$).

\paragraph{What it changes.}
The weak association of the earlier criterion (30 of 55 against 7 of 67) came from its truncation; with the full
identity the stall diagnosis gives 30 of 38 predicted-dead runs dead and 7 of 126 others. HT routes with $k=5$ never
stall (0 of 144) because HT leaves need $k\ge6$; the drop of the stalls with spare channels is the protection of (6),
while their rise on capped routes from small initialization is reproduced per run but its mechanism is not isolated.
Steep MPS routes: if pinned routes escape as conjectured in (5), a pinned end gives the exponent $2k-6$ and end
amplitudes $\propto\isc^{k-3}$, and a partly decayed end values in between, qualitatively as measured ($k=5$: $3.3$--$3.8$;
not tested per seed). Prior work: a conserved quantity that decides which side of a unit can vanish is Lemma~B.1 of
\citet{kunin2025alternating}; \citet{furman2026benign}, Thm~G.1, give a spurious minimum away from the minimum-norm
set, and the trapped endpoints here are reached because the conserved imbalance, of order $\isc^2$, is as large as the
dormant entries.

\subsubsection{Provenance and corrections after review}\label{app:supp-s2s-provenance}
The numbers of \cref{sec:s2s} come from the committed result tables of the experiments \path{s2s_core},
\path{locality}, \path{order} and \path{regime_map} (\cref{app:reproducibility}), checked against them in two
verification reviews per experiment and a final review of the claims. Three numbers of \cref{sec:s2s-order} exist only in
the independent recomputation of the order experiment (scripts in \path{notes/reviews/scripts/verify_order/}): the fits
over the three smallest values of $\isc$ (\path{v_sigma.py}, \path{v_sigma2.py}, \path{v_sigma3.py}), the paired McNemar
comparison of each run's order with its own LO race (\path{v_multi.py}), and the margin of the stall rule of the
order runs, which had no false positive in 22 runs without any stall stop (\path{v_multi.py}). The same review
reproduced $\mathcal{T}(Z)$ for the 160 binary4 and MPS4 first-escape jobs with an independent integration of the spherical
flow, to $2\cdot10^{-10}$. The independent cross-check of the locality results is the pilot
\path{experiments/pilot_h2_locality/} (its own code). Corrections after review: the norm law
$\tev1\propto\norm{\params_0}^{-(\Kc-2)}$, the first guess of the research plan, is only a lower bound
(\cref{sec:s2s-escape}); the binary4 sequence, its $\isc^{-1}$ plateaus and its strict wrong-order saddle are known for
deep linear networks and are reported as checks; the single-route consistency of the bond spectra is a check of the
Steiner edges, not a prediction success rate; the law $\isc^{-(2k-6)}$ for steep MPS routes was the pilot's conjecture;
the claims withdrawn in the reviews are not restated.


